% related_work -- owner: TBD
\section{Related Work}

Conversational pushback studies examine whether models revise their answers when challenged and how those revisions affect accuracy (Laban et al., 2023; Xie et al., 2024; Sinha, 2026). Closer to our setting, user influence on analytical conclusions has already been shown, so it is not our contribution. Allen and Peterson (2026), the most similar design, vary desired-outcome framing in requests to analyze synthetic data with code execution and assess conclusions against known ground truth; the desired outcome arrives in the first request, and responses are coded for the estimate rather than the strength of the conclusion. Balani and Panda (2026) hold statistical summaries fixed across framing conditions and use an LLM judge to separate factual departures from changes in tone, but give the model no analysis to run. Asher et al.\ (2026) give two coding agents the replication data of four published null results and vary hypothesis framing and pressure for significance; estimates stayed stable until a prompt reframed specification search as uncertainty reporting, but published data carry no known ground truth.

Huang and Zhang (2026) replay the same item and committed answer with and without a pressure statement, and Ping et al.\ (2026) compare false claims introduced in an initial medical question with challenges delivered after the model has answered; both are precedents for controlled replay and for comparing when pressure arrives. Bertran et al.\ (2026) find that analyst personas change conclusions even among runs that pass an AI methodological audit, which shows why a changed conclusion alone does not establish an error.

Our outcome comes from the literature on spin: reporting that presents a result as more favorable than the evidence supports (Boutron et al., 2010; Lazarus et al., 2015), which shifts how expert readers judge a treatment's benefit (Boutron et al., 2014) and which LLMs are susceptible to and carry into their summaries (Yun et al., 2025; Peters \& Chin-Yee, 2025). These studies take the spin as given in the source; we ask whether a user's expectation produces it in an analysis the model ran itself. We generate the data, so the correct conclusion is known to us but hidden from the user and the model, and the model runs its own causal analysis, which can fork in the specifications it chooses (invisible forking; Ashton, 2026a, \S\,2.3). Here the fork we control is the user's next turn: a one-sentence expectation, stated after the analysis or before it. We do not test the mechanism behind a shift or its effect on readers.

